% Ejemplo: malla de un área CON estándares y DBA (columnas originales del formato del colegio).
% Contenido ilustrativo de Matemáticas 6.°, periodo I (no es una malla oficial; los estándares y
% DBA deben copiarse literalmente de los documentos del MEN al hacer una malla real).
\documentclass[malla]{planeacion}
\graphicspath{{../}}
\materia{Matemáticas}
\grado{6.°} \periodo{I} \anio{2027} \docente{Docente del área}
\begin{document}
\begin{malla}{Pensamiento numérico y sistemas numéricos: resuelvo y formulo problemas en contextos de medidas relativas y de variaciones en las medidas; utilizo números racionales en sus distintas expresiones para resolver problemas.}
  \fila{\textbf{1. Números naturales.} Sistema de numeración decimal, valor posicional, orden y
    operaciones; propiedades de la adición y la multiplicación.}
    {Pensamiento numérico}{Resolución de problemas con operaciones combinadas.}
    {\punto Razonamiento \punto Resolución de problemas}
    {\punto Interpreta los números naturales y sus operaciones en contextos cotidianos.}
  \fila{\textbf{2. Potenciación, radicación y logaritmación.} Relación entre las tres operaciones;
    cuadrados y cubos perfectos.}
    {Pensamiento numérico}{Relaciones entre operaciones inversas.}
    {\punto Comunicación \punto Razonamiento}
    {\punto Utiliza las propiedades de las operaciones para simplificar cálculos.}
  \fila{\textbf{3. Teoría de números.} Múltiplos y divisores, números primos y compuestos,
    descomposición en factores primos, m.c.m. y m.c.d.}
    {Pensamiento numérico}{Uso de la descomposición factorial en problemas de reparto.}
    {\punto Resolución de problemas \punto Modelación}
    {\punto Resuelve problemas de reparto y de periodicidad con m.c.m. y m.c.d.}
  \fila{\textbf{4. Fracciones.} Fracción como parte, como cociente y como razón; fracciones
    equivalentes; orden.}
    {Pensamiento numérico}{Representaciones de los números racionales.}
    {\punto Comunicación \punto Modelación}
    {\punto Representa fracciones en la recta numérica y con diagramas, y las compara.}
  \fila{\textbf{5. Datos.} Tablas de frecuencia, diagramas de barras y pictogramas; moda.}
    {Pensamiento aleatorio}{Recolección y organización de datos del entorno escolar.}
    {\punto Comunicación \punto Razonamiento}
    {\punto Organiza datos en tablas y gráficos y los interpreta para responder preguntas.}
\end{malla}
\end{document}
